\section{Anharmonic, disordered and quantum solids}\label{app:E.19}

Solid helium near melting has a Lindemann ratio of about a quarter: the zero-point amplitude of an atom is about a quarter of the distance to its neighbor (\cref{sec:11.5}), and the harmonic lattice is not a valid starting point. Strontium titanate and potassium tantalate have soft modes whose frequencies depend on temperature at order one; lithium hydride has a light atom whose zero-point motion is large. For these targets the multiphonon response of \cref{sec:11.5} has no small parameter. Glasses lack the lattice that the response tables assume, and a random medium can scatter a strongly interacting particle coherently over its whole volume.

\subsection{Strongly anharmonic and quantum crystals}

Every integer frequency moment of the self response is an equal-time correlation function. The moments $\mathscr{M}^k$, $k = 0,\dots,4$, of the self response of \cref{sec:8.3} are therefore exact from a static path-integral simulation, for any anharmonicity, any temperature and any Lindemann ratio. This holds in the Born--Oppenheimer approximation, where the nuclei move on a smooth energy surface and every moment is finite; with point nuclei and the full Coulomb interaction the fourth moment diverges (\cref{sec:8.3}). The Markov--Krein band of \cref{sec:11.5} then brackets every response with these five moments; it bounds bin-integrated rates in the bulk of the response, and threshold rates only from above. The calculation therefore has two parts, exact moments and an approximate real-time response tested against them.

For KTaO$_3$, SrTiO$_3$ and LiH we want $\mathscr{M}^0$--$\mathscr{M}^4$ of the self response at $|\vect{q}| = 0.5$ to \SI{20}{keV}, the band per $\vect{q}$ at the millikelvin operating temperature, and the self response $S_{\mathrm{self}}(\vect{q},\omega)$ (per volume, in \si{eV^2}, the units of \cref{sec:4.2}) at 50 and \SI{100}{K}, where the real-time methods below hold. On the truncation label $\mathfrak{T}$ (\cref{sec:8.5}) the moments are exact in $n_A$; the dynamics is approximate at all orders of $n_A$, and the moment test measures its error. The static part is path-integral molecular dynamics in i-PI~\cite{kapil2019ipi} with a machine-learned potential (\cref{app:E.11}) trained on PBE and validated on displaced configurations to three standard deviations of the path-integral distribution. The cells are cubic $n\times n\times n$ supercells with $n = 4$, 6 and 8 (320, 1080 and 2560 atoms) for the perovskites, even $n$ so that the soft mode at the R point is commensurate, and $4^3$ to $6^3$ conventional cells (512 to 1728 atoms) for LiH. The number of beads $P$ satisfies $\beta\omega_{\mathrm{max}}/P \lesssim 0.3$; for SrTiO$_3$ at \SI{50}{K} with $\omega_{\mathrm{max}} \approx \SI{100}{meV}$ that is $P \approx 80$, doubled until $\mathscr{M}^2$ changes by less than \SI{1}{\percent}. Colored-noise thermostats may accelerate the static moments but must not be used in the dynamics, which they alter.

The dynamic part compares three approximate methods. Thermostatted ring-polymer molecular dynamics damps only the internal ring-polymer modes, with friction $\gamma_k = \lambda\omega_k$ at $\lambda = 1/2$ and checks at $1/4$ and 1, which preserves the exact limits of the method~\cite{rossi2014remove}; the centroid is not thermostatted during the dynamics. Centroid molecular dynamics, with the time step reduced threefold, is the second method, and their difference is the method uncertainty. The third is the stochastic self-consistent harmonic approximation~\cite{monacelli2021stochastic}, first static, then with the dynamical bubble, then the full time-dependent form~\cite{monacelli2021time}; the bubble-to-full difference measures the order. The time step is at most $0.1/\omega_P$ of the stiffest free ring-polymer mode, typically 0.1 to \SI{0.25}{fs}, with centroid energy drift below $10^{-4}$ per ps. Resolving $\delta\omega = \SI{0.5}{meV}$ needs trajectories of $20\pi/\delta\omega \approx \SI{80}{ps}$ at $q \lesssim q_{\mathrm{BZ}}$, shorter at large $q$ where the correlation decays; 100 independent initial conditions give \SI{1}{\percent} statistical error on bin integrals. Every approximate response must lie inside the band of its exact moments, and the response must change by less than \SI{2}{\percent} from one cell size to the next while the finite-size gap $2\pi c_\nu/L$ stays below the lowest frequency bin. A response outside the band is flagged, not used.

Solid $^4$He is computed the same way with the Aziz pair potential, which needs no fitting. The simulation is path-integral molecular dynamics at a molar volume of \SI{21}{cm^3.mol^{-1}} and \SI{1}{K}, 180 to 432 atoms and $P = 64\to128$ beads (the Debye energy is \SI{2}{meV}, so $\beta\omega_{\mathrm{D}} \approx 25$), with an open-path estimator for the momentum distribution. The band is computed for the self and the coherent response; the coherent moments are the static structure factor, the $f$-sum, the current correlation and the third-moment correlator, all equal-time. Bose exchange is negligible for the statics at this temperature but is checked once by permutation sampling. At 1 to \SI{40}{GPa}, where the Lindemann ratio drops to about 0.1, the self-consistent harmonic force constants from a potential trained on dispersion-corrected PBE replace the static Hessian in the phonon response of \cref{sec:11.5}, converged to \SI{2}{\percent} in the frequencies and \SI{3}{\percent} in the displacement correlations.

Ring-polymer and centroid methods give Kubo-transformed correlations; converting to $S$ multiplies by $\beta\omega/(1 - \ee^{-\beta\omega})$, which amplifies their error by $\beta\omega$, so they are usable only for $T \gtrsim \omega/10$ and not at a \SI{10}{mK} operating temperature. Their short-time error for the nonlinear operator $\ee^{\ii\vect{q}\cdot\vect{u}}$ is already $O(t^3)$. The time-dependent self-consistent method excites at most two phonons at a time, so at $q \gg q_{\mathrm{BZ}}$ it can supply only the dressed one-phonon propagator inside the Gaussian form. Tunneling splittings in double wells are missed by centroid-like methods. At millikelvin, therefore, the response of these targets is bounded by the band alone. The ring-polymer runs cost $10^3$ to $10^4$ core-hours per material and temperature, and centroid dynamics three times more. The self-consistent harmonic runs cost about $10^2$ core-hours per material, path-integral helium at melting about $10^3$, and high-pressure helium about $10^4$ per pressure. The memory is a few gigabytes. The moments and bands enter \cref{sec:11.5}.

\subsection{Weak anharmonicity}

In silicon the anharmonicity is small, and the question is the size of its first correction to the multiphonon response above the one-phonon band, where the low-mass reach is set. Time reversal fixes the structure. For a target invariant under time reversal $S_{\kappa\kappa}(\vect{q},\omega) = S_{\kappa\kappa}(-\vect{q},\omega)$, and a term of first order in the cubic force constants $\Phi^{(3)}$ is of third degree in $\vect{q}$, so it cancels. The leading correction is therefore of second order in $\Phi^{(3)}$ or first order in $\Phi^{(4)}$, the level $n_A = 2$ of \cref{sec:8.5}: the one-phonon propagator dressed by the bubble and loop self-energies, the connected four-point function, and the square of the connected three-point function. We want the relative correction $\delta_{\mathrm{anh}}(\vect{q},\omega) = S/S^{\mathrm{harm}} - 1$ of \cref{sec:11.5}, for silicon first, then germanium and gallium arsenide.

The cubic force constants come from finite displacements in a $4\times4\times4$ primitive supercell (128 atoms) with a pair cutoff of \SI{5}{\angstrom}, about 300 supercell DFT runs, through phono3py~\cite{togo2015distributions}; the quartic ones from hiPhive~\cite{eriksson2019hiphive} or TDEP. The bubble self-energy gives the full spectral function on a $24^3$ grid, not on-shell linewidths, and the quartic loop gives the shift. The three- and four-point functions are mode sums over four branch labels with contour kernels built from the thermal propagator, about $10^3$ core-hours. The dense cubic constants of the 128-atom cell have $384^3$ entries, \SI{0.45}{GB} in double precision; the quartic set within the cutoff is smaller. The spread under the supercell from $3^3$ to $4^3$, the cutoff from 4 to \SI{5}{\angstrom} and the grid from $24^3$ to $32^3$ is the uncertainty of the correction. The first-order cubic term must vanish numerically, which checks the placement of the time arguments on the contour, and the correction must satisfy the exact linear constraints of \cref{sec:11.5}, anchored by the kinetic energy of the struck atom from the path-integral calculation of \cref{app:E.17}. The total cost is $10^3$ to $10^4$ core-hours.

\subsection{Glass models}

Amorphous SiO$_2$, silicon and silicon nitride need structures before they need responses. The models are made by melt-quench molecular dynamics with a machine-learned potential (\cref{app:E.11}), 648 to 5184 atoms and at least five independent quenches, validated on the measured density and pair distribution. From each model come the full Hessian and the Born charges $Z^*_I$ of every atom, either by DFPT on 192- to 648-atom subcells or from a charge-equilibration model fitted to them. The Hessian of 5184 atoms has $15552^2$ entries, \SI{1.9}{GB} in double precision, and its dense diagonalization costs about $4\times10^{12}$ floating-point operations.

The infrared coefficient $C^{\mathrm{IR}}(\omega) = A^{\mathrm{IR}} + B^{\mathrm{IR}}\omega^2$ of glass absorption below the Ioffe--Regel frequency, \cref{eq:16.7-CIR,eq:16.7-AB}, is the first quantity that enters \cref{sec:16.7}. $A^{\mathrm{IR}}$ and $B^{\mathrm{IR}}$ follow from the Born-charge structure factor $S^Z_\nu(k)$ of static configurations and a charge model, without the Hessian, on wavevectors $k \le \SI{0.3}{\angstrom^{-1}}$ of the periodic cell; the Hessian gives the full $C^{\mathrm{IR}}(\omega)$ as a cross-check up to the boson peak. The second quantity is the distinct part of the response at two and three phonons. It needs the configurational averages, at fixed pair distance $r$, of the longitudinal and transverse parts of the pair displacement tensor, its longitudinal--transverse vector and its traceless transverse part, binned in $\Delta r = \SI{0.02}{\angstrom}$ up to \SI{15}{\angstrom}; the third is the site-resolved phonon density of states. The ensemble is validated against the measured terahertz absorption of the same glass~\cite{taraskin2006universal}.

The pair distribution, the moments, the density of states and $C^{\mathrm{IR}}$ change by less than \SI{5}{\percent} between 1728 and 5184 atoms; the box is at least \SI{4}{nm}, so that $2\pi/L$ lies below the correlation length of the charges; and $\sum_I Z^*_I = 0$ in every configuration. The 5184-atom box of SiO$_2$ is \SI{4.3}{nm} across; the 1728-atom box, \SI{3.0}{nm}, serves only as the smaller member of the size test. A locally neutral charge model forces $A^{\mathrm{IR}} = 0$ by construction and must not be used; the model has to reproduce the site-charge fluctuations of the subcell DFPT. The finite-size gap $2\pi c_\nu/L$, 3 to \SI{4}{meV} at \SI{4}{nm}, removes the lowest modes, and the dependence on quench rate is a model systematic. The cost is $10^4$ core-hours for the molecular dynamics, $10^3$ for the subcell Born charges and $10^2$ per model afterward.

\subsection{Coherent force from a random medium}

A dark-matter particle with a large cross-section on nuclei sees a solid as the optical potential $V_{\mathrm{opt}}$ of \cref{eq:16.7-Vopt}, and the Born treatment of scattering off density fluctuations holds only for $\sigma_{\chi N} \ll \sigma_{\mathrm{crit}}$ (\cref{eq:16.7-scrit}). Above it the force on a test mass comes from reflection at its faces. A plane wave with normal velocity $v_n$ meets the step with reflection probability $R(v_n) = |(k_n - k')/(k_n + k')|^2$, $k'^2 = k_n^2 - 2\mu_{\chi N}V_{\mathrm{opt}}$; a slab many de Broglie wavelengths thick reflects, averaged over its thickness, $R_{\mathrm{slab}} = 2R/(1 + R)$. Each reflected particle returns momentum $2m_\chi v_n$, and the flux on a face of outward normal $\hat{\vect{n}}$ is $\rho_\chi/m_\chi$ times the inward normal velocity, so the pressure is
\begin{equation}
  P_{\mathrm{wall}}(\hat{\vect{n}}) = 2\rho_\chi\int\dd^3v\,f_\chi(\vect{v})\,(\vect{v}\cdot\hat{\vect{n}})^2R_{\mathrm{slab}}(\vect{v}\cdot\hat{\vect{n}})\,\Theta(-\vect{v}\cdot\hat{\vect{n}}) ,
  \label{eq:E.19-wall}
\end{equation}
while transmitted particles leave with their initial velocity and transfer nothing on average. What the equivalence-principle test measures is the differential acceleration $\Delta a(\sigma_{\chi N},m_\chi)$, in \si{m.s^{-2}}, of the two cylinders of the MICROSCOPE equivalence-principle test (a platinum alloy and Ti--6Al--4V), summed over faces and averaged over the standard halo with the Earth's velocity. The scale is set by $\rho_\chi v^2 \approx \SI{5e-11}{Pa}$.

The calculation evaluates \cref{eq:E.19-wall} per face of the axisymmetric geometry, with one-dimensional transfer matrices along the local normal (the curvature radius is far larger than the de Broglie wavelength, \SI{60}{\micro m} at \SI{20}{eV}), and a four-dimensional Monte Carlo over velocity and face. The code is about 300 lines and runs in CPU minutes. Below $\sigma_{\mathrm{crit}}$ the same code evaluates the distorted-wave Born form, with the density correlation $\langle\delta\nd\,\delta\nd\rangle$ taken from small-angle X-ray scattering on the actual alloy. As $V_{\mathrm{opt}}/E_{\mathrm{kin}} \to 0$ the result must reproduce the Born random-medium rate of \cref{sec:16.7}. Roughness on the de Broglie scale makes the reflection diffuse, and the nested cylinders reflect particles back and forth; both are included in the Monte Carlo.

\subsection{Chiral phonons in molecular frameworks}

A chiral phonon carries angular momentum $L^{\mathrm{ph}}_{\vect{Q}\nu} = -\ii\sum_\kappa\vect{e}^*_{\kappa\nu}(\vect{Q})\times\vect{e}_{\kappa\nu}(\vect{Q})$. The proposal of \cref{sec:14.5} reads out the phonons that end, after anharmonic decay, in a long-lived chiral band with $|L^{\mathrm{ph}}| > L_{\mathrm{th}}$~\cite{romao2023chiral}. Its yield is the mean number of such phonons per phonon created in mode $\lambda = (\vect{Q}\nu)$, averaged over the created spectrum. Although the decay branches, the expectations obey a linear system. Let $\mathsf{M}_{\lambda\lambda'}$ be the mean number of daughters in $\lambda'$ per decay of $\lambda$: the three-phonon rates into unordered pairs, counted twice for identical daughters, plus the elastic isotope rate, all divided by the total rate of $\lambda$, which includes boundary loss and sensor capture. Let $b_\lambda = 1$ for modes in the terminal chiral set, those with $|L^{\mathrm{ph}}| > L_{\mathrm{th}}$ and a decay time longer than the observation time $t_{\mathrm{obs}}$, and $b_\lambda = 0$ otherwise, with the rows of $\mathsf{M}$ zero for terminal modes. Conditioning on the first decay gives $\mathcal{N}^{\mathrm{ch}}_\lambda = b_\lambda + \sum_{\lambda'}\mathsf{M}_{\lambda\lambda'}\mathcal{N}^{\mathrm{ch}}_{\lambda'}$, so
\begin{equation}
  \boldsymbol{\mathcal{N}}^{\mathrm{ch}} = (\mathbb{1} - \mathsf{M})^{-1}\vect{b} .
  \label{eq:E.19-branch}
\end{equation}
Without isotope scattering $\mathsf{M}$ is nilpotent in energy order. With it, $\mathbb{1} - \mathsf{M}$ is still invertible as long as every closed set of elastically connected modes has some decay or loss.

The first framework is InF$_3$(4,4$'$-bipyridine), whose chiral phonon bands Romao et al. computed~\cite{romao2023chiral}. The rates come from phono3py on the framework cell, at $n_A = 2$ (Fermi's golden rule with two cubic vertices) with harmonic $L^{\mathrm{ph}}$, on a $10\times10\times20$ grid; \cref{eq:E.19-branch} is a sparse solve. Cubic force constants from DFT cost about $10^5$ core-hours for a metal--organic framework; a machine-learned potential validated on 50 DFT supercells brings this to $10^3$. The grid, the cutoff of the cubic constants and the smearing of energy conservation are converged, and the sensitivity to $t_{\mathrm{obs}}$ is reported. A direct Monte Carlo of the branching process with the same rates must reproduce $\boldsymbol{\mathcal{N}}^{\mathrm{ch}}$ within its statistical error, and the energy in the terminal set plus the energy lost must equal the energy created. $L^{\mathrm{ph}}$ is not unique inside a degenerate subspace, so it is diagonalized there. It is not conserved in decay, so the label is recomputed for each daughter. And the code must state whether its rates are per ordered or per unordered pair, since the factor two moves the answer at order one.

Frameworks with lanthanide ions, whose host and ion \cref{sec:14.5} selects, add one-phonon amplitudes through the crystal field. Let $P_{\mathrm{K}}(\vect{u})$ be the projector on the ground Kramers doublet, of energy $E_{\mathrm{K}}$, of the crystal-field Hamiltonian $H_{\mathrm{CF}}(\vect{u})$ of the $4f$ shell at displacement $\vect{u}$. Let $O$ be the spin $\vect{s}_e$ or the orbital moment $\vect{L}$, and $R = \sum_{c\notin\mathrm{K}}|c\rangle\langle c|/(E_{\mathrm{K}} - E_c)$. First-order perturbation theory gives $\partial P_{\mathrm{K}} = R\,\partial H_{\mathrm{CF}}\,P_{\mathrm{K}} + P_{\mathrm{K}}\,\partial H_{\mathrm{CF}}\,R$, and the amplitude follows from $\partial(P_{\mathrm{K}}OP_{\mathrm{K}})$ contracted with the phonon eigenvector, $\vect{e}_{\kappa\nu}(\vect{Q})/\sqrt{2NM_\kappa\omega_{\vect{Q}\nu}}$. The projector form avoids derivatives of individual doublet states, which depend on an arbitrary phase. The inputs are the derivatives of the 27 Stevens parameters ($k = 2,4,6$) and of the $g$ tensor with respect to symmetry-adapted displacements of an embedded cluster (the ion, its first coordination shell and whole ligands, in Ewald point charges with model potentials on boundary cations), computed by central differences of $\pm\SI{0.01}{\angstrom}$. The electronic structure is a state-averaged complete-active-space self-consistent field (CASSCF) calculation over the ground term with spin--orbit mixing, a perturbative correction of the crystal-field splittings, a scalar-relativistic Hamiltonian and ANO-type basis sets. The derivatives must be linear over 0.005 to \SI{0.02}{\angstrom} to \SI{2}{\percent} and stable to \SI{10}{\percent} under enlargement of the active space, with the cluster extended to its second shell and the basis enlarged as further checks. The roots must be tracked by overlap, because they can exchange order under displacement, and $\vect{s}_e$ and $\vect{L}$ are kept separate since dark matter couples to them differently. Each geometry costs 0.5 to 2 core-hours, $10^2$ to $10^3$ per inequivalent site. Both results enter \cref{sec:14.5}.
